% 工程内兼容修复；外部 Template/ 与 Styles/ 保持只读。
\NewCommandCopy\BookOriginalNewCJKfontfamily\newCJKfontfamily
\RenewDocumentCommand{\newCJKfontfamily}{m O{} m O{}}{%
  \ifstrequal{#3}{FZCuSong-B09S}{%
    \IfFontExistsTF{FZCuSong-B09S}{%
      \BookOriginalNewCJKfontfamily#1[#2,#4]{#3}%
    }{\BookOriginalNewCJKfontfamily#1[BoldFont=SimSun,AutoFakeBold=true]{SimSun}}%
  }{\BookOriginalNewCJKfontfamily#1[#2,#4]{#3}}%
}
\usepackage{../../Styles/mystyle-elegantbook}
\RenewCommandCopy\newCJKfontfamily\BookOriginalNewCJKfontfamily
% 仿宋常规字形已有，补用同套黑体的粗体，避免证明内局部定义的粗体替代警告。
\newCJKfontfamily\BookFangSong[BoldFont=FandolHei-Bold.otf]{FandolFang-Regular.otf}
\renewcommand{\cfs}{\BookFangSong}
% 分款引用的可见编号来自真实计数器，链接指向最深层分款。
\renewcommand{\subcref}[2]{\hyperref[#2]{\namecref{#1}~\ref*{#1}\ref*{#2}}}
\renewcommand{\subsubcref}[3]{\hyperref[#3]{\namecref{#1}~\ref*{#1}\ref*{#2}\ref*{#3}}}
\crefname{chapter}{章}{章}
\crefname{equation}{公式}{公式}
\crefname{exer}{练习}{练习}
\crefname{exam}{例题}{例题}
\crefname{axiom}{公理}{公理}
\crefname{tcb@cnt@axiom}{公理}{公理}
\elegantnewtheorem{principle}{原理}{thmstyle}{principle}
\crefname{principle}{原理}{原理}
\crefname{tcb@cnt@principle}{原理}{原理}
\setlist[enumerate,1]{label=(\alph*),ref=(\alph*)}
\setlist[enumerate,2]{label=(\roman*),ref=(\roman*)}
\DeclareMathOperator{\card}{card}

% 对中文正文统一多重引用的连接词；编号仍由 cleveref 动态生成。
\providecommand{\crefrangeconjunction}{}
\renewcommand{\crefrangeconjunction}{至}
\providecommand{\crefpairconjunction}{}
\renewcommand{\crefpairconjunction}{和}
\providecommand{\crefmiddleconjunction}{}
\renewcommand{\crefmiddleconjunction}{、}
\providecommand{\creflastconjunction}{}
\renewcommand{\creflastconjunction}{和}
% 预留证明标题及其首行的空间，防止标题单独留在页底。
\usepackage{needspace}
\BeforeBeginEnvironment{proof}{\Needspace{5\baselineskip}}
\BeforeBeginEnvironment{exercise}{\Needspace{9\baselineskip}}
\BeforeBeginEnvironment{remark}{\par}
\AfterEndEnvironment{remark}{\par}
% 原模板仅有编号 example；局部例子沿用其字形和颜色，不递增计数器。
\NewDocumentEnvironment{example*}{O{}}{%
  \par\noindent{\normalfont\songti\bfseries\color{main}\examplename\ #1 }\rmfamily
}{\par\ignorespacesafterend}

\crefname{enumi}{项}{项}
\crefname{enumii}{项}{项}
